\chapter{Physical Evidence}
\label{chap:physical-evidence}

\classification{Evidence status: physical-evidence inventory. Confidence baseline: High for documented existence of major exhibits where source-mapped to DOC-0001; Moderate or Low where item detail, preservation, current custody, or OCR transcription remains incomplete. This chapter does not infer ownership, motive, cause of death, or significance beyond documented evidence.}

\section{Overview of Physical Evidence}

The physical-evidence layer contains 34 catalogue entries, 14 scene-reconstruction components, 34 chain-of-custody confidence entries, 10 preservation-history entries, and 10 unresolved physical-evidence questions. The chapter uses those records as an inventory system rather than as a theory of the case.

The catalogue covers the body and scene context, body clothing, pocket property, transport tickets, suitcase and suitcase contents, Tamam Shud material, Rubaiyat material, fingerprints, dental records, photographs, the death mask or bust, biological specimens, environmental observations, thread, sand, vegetation, and later preservation records. The strongest entries are those supported by direct inquest testimony and page-mapped documentary references. The weakest entries are those known only through OCR-dependent inventory wording, later reproduction, or preservation leads without original-custodian confirmation.

Custody is the principal limitation. No physical-evidence entry is classified as having a complete custody trail. The custody matrix records 29 entries as fragmentary, 4 as unknown, and 1 as not applicable. This does not negate the existence of the exhibits; it limits the weight that later examination or interpretation can safely place on them.

\begin{table}[htbp]
\centering
\caption{Physical evidence category summary}
\label{tab:physical-evidence-category-summary}
\begin{tabularx}{\textwidth}{>{\raggedright\arraybackslash}p{0.34\textwidth}r>{\raggedright\arraybackslash}X}
\toprule
Category & Count & Inventory use \\
\midrule
Physical evidence catalogue entries & 34 & Complete working inventory for object-level discussion. \\
Scene reconstruction components & 14 & Structured reconstruction of location, body position, scene observations, and evidence placement. \\
Chain-of-custody entries & 34 & Custody confidence and gap register for each catalogue item. \\
Preservation-history entries & 10 & Current preservation, local reproduction, and access-status register. \\
Unresolved physical-evidence questions & 10 & Priority acquisition and verification targets. \\
\bottomrule
\end{tabularx}
\end{table}

\section{Discovery Scene}

The scene reconstruction places the body on the Somerton Beach foreshore near the seawall opposite the Somerton Crippled Children's Home. The body-location component records the body on its back, with head and shoulders supported by the seawall and the head inclined to the right. The orientation record describes the feet toward the sea or west, the right arm doubled over with palm upward and fingers bent, and the left arm alongside the body.

The scene record identifies a partly smoked cigarette near the right cheek, right lapel, or collar area. It also records that the evening witness did not see a cigarette in the raised hand. The cigarette is therefore catalogued as a scene object and body-adjacent observation, not as a recovered and tested exhibit.

Distance and environmental information remain incomplete. One evening witness estimated a viewing distance of 15 to 20 yards and the location as within about a yard of steps. The database records the beach as an open public location in a summer context. It also records warm or hazy conditions from testimony, no undue sand disturbance noted by Moss, and no original scene sketch acquired.

Scene documentation is limited. The project has source references to body photographs and later exhibit photographs, but the original scene-photograph set is not preserved locally as publication-cleared imagery. The absence of a scene grid, precise measurements, original scene diagram, and rights-cleared image set means the reconstruction must remain component-based rather than diagrammatic.

\section{Clothing}

The body-clothing ensemble is catalogued as PE-0003 and maps to EV-0005. It includes coat, shoes, shirt, pullover, jockey underpants, singlet, socks, trousers, and one tie. The catalogue records the clothing as fairly good quality, with the shoes described as very clean or practically new in the testimony layer. The current location and present survival status of the clothing are unknown.

The tie evidence is split across body and suitcase contexts. PE-0004 records the tie worn by the deceased and two ties in the suitcase, including one bearing the name T. Keane. The catalogue does not resolve which tie carried which mark, whether any tie survives, or whether the marked name establishes ownership.

The shoes are catalogued separately as PE-0005. The record supports existence and condition observations, including clean or polished condition and limited sand observations. It does not support inference about gait, occupation, travel pattern, or social status beyond the recorded descriptions.

The socks are catalogued as PE-0006. Cleland's observation of barley-grass seed material is also recorded in the vegetation entry. The sock record has moderate confidence because exact sock type, condition, retention status, and whether any vegetation was preserved remain unknown.

Removed or missing clothing tags and labels are catalogued as PE-0007. The record supports that tags or labels were missing or removed from some body and suitcase articles. It also preserves the important limitation that Cleland did not find threads proving recent removal. The chapter therefore records label removal as an observed condition, not as proof of timing, intent, or concealment.

\section{Personal Property}

The personal-property record begins with items found in the initial clothing search. The catalogue records the Henley Beach railway ticket as PE-0008 and the Municipal Tramways Trust bus ticket as PE-0009. Both are strong transport-property entries for issue and recovery context, but neither identifies the ticket purchaser or proves the full route travelled.

The cigarette-related property is separated into two records. PE-0002 is the partly smoked cigarette observed near the face or collar area. PE-0010 is the Army Club cigarette packet containing Kensitas cigarettes. The catalogue preserves the distinction between observation of the partly smoked cigarette and inventory of the cigarette packet. It records no confirmed saliva, fingerprint, brand-comparison, or toxicological examination of the cigarette.

PE-0011 records a Bryant and May matches box. PE-0012 records two combs or a metal comb reported in body property. PE-0013 records a Juicy Fruit chewing gum packet. These entries are inventory items only. The catalogue does not support behavioural inference from them.

The catalogue also records PE-0033, a low-confidence cash entry of 2s 6d reportedly found in a trouser pocket according to OCR-dependent Leane text. This item requires visual confirmation before use as a narrative detail. The record is retained in the inventory because it may be part of the body-property layer, but it is not treated as fully established.

\section{Suitcase and Contents}

The suitcase is catalogued as PE-0017 and maps to EV-0003. It was found unclaimed at the Adelaide Railway Station luggage or cloak room. The catalogue records that Leane located it in unclaimed luggage on 14 January 1949 and that police took possession on 19 January 1949. The corresponding luggage ticket is PE-0018; North produced the ticket, and the exact depositor identity remains unknown.

The suitcase record supports documented existence, police possession, and internal association evidence. It does not prove ownership. It also does not establish a complete custody trail, a complete inventory, or current physical survival. The preservation register records the suitcase's current status as unknown, with reported destruction not primary-verified in this repository.

Suitcase clothing and textiles are catalogued as PE-0019. The record includes dressing gown or cord, laundry bag, singlets, underpants, ties, slippers, trousers, sports coat, shirt, pajamas, scarf, towel, and related textiles. The item-level inventory is moderate confidence because the record depends on OCR-distorted testimony and lacks a clean verified layout.

Suitcase tools and personal articles are catalogued as PE-0020. The record includes scissors, knife, stencil brush, razor strop, lighter, razor, shaving brush, screwdriver, pencils, toothbrush or paste, dish, soap dish, hair pin, safety pins, collar studs, button, spoon, broken scissors, boot polish, and air-mail stickers. The catalogue records limited examination of the stencil-related articles, but does not infer occupation from them.

Pencils are separately catalogued as PE-0021. The record notes six pencils, most Royal Sovereign, with OCR indicating three H-type or drafting pencils. This supports the presence of writing implements in the suitcase inventory. It does not establish a connection to the Rubaiyat code or to a technical occupation.

The warm sepia, orange, or tan thread is catalogued as PE-0022. The record supports testimony that thread found in the suitcase corresponded microscopically with thread used on body-clothing repairs. It is association evidence, not ownership proof by itself, and the physical retention status of thread samples remains unknown.

\section{Tamam Shud Slip}

The Tamam Shud paper slip is catalogued as PE-0014 and maps to EV-0001. It was found in a concealed or difficult-to-locate trouser fob pocket after the initial search had not located it. The record supports existence, fob-pocket association, paper-type inquiry, and later exhibit handling.

The slip's custody is fragmentary. The exact date and time of discovery are unresolved. The finder sequence is unresolved. Brown received the slip from Leane for inquiries, and Cleland described the pocket as easy to miss, but the catalogue does not establish a complete physical custody chain from discovery to present.

The known examination layer is limited to textual and paper comparison, bookshop or library inquiry, and exhibit handling. No fibre, chemical, or modern physical examination report is acquired in the repository. The current physical status and custodian are unknown.

\section{Rubaiyat}

The Rubaiyat copy or reference-book material is catalogued as PE-0015 and maps to EV-0002. The record supports that Brown compared the slip with bookshop and library copies and that the coroner associated the paper with a second edition of Fitzgerald's translation of the Rubaiyat of Omar Khayyam.

The associated original book remains a custody problem. Exact recovery, owner or finder, edition details, tear-match status, and current location are unresolved in the acquired primary-source layer. The chapter therefore records the Rubaiyat as a documented comparison and association problem, not as a settled ownership exhibit.

The code or markings image is catalogued separately as PE-0016. A digital reproduction is preserved locally as an Abbott archive image, while original book or police custody remains unresolved. This chapter does not analyze the code. Documentary, cryptographic, and handwriting issues belong to the later document-analysis and intelligence-assessment units.

\section{Biological Evidence}

Fingerprint evidence is catalogued as PE-0025 and maps to EV-0015. The fingerprint chart is preserved locally as a digital reproduction, while Durham's testimony records fingerprinting and circulation to Australian states, New Zealand, and important overseas bureaus. The current evidentiary value is high for existence and moderate for custody because the original chart and complete bureau-response logs have not been acquired.

Dental evidence is catalogued as PE-0026 and maps to EV-0016. Dwyer prepared a dental chart, and a separate chart reproduction is preserved locally. The record is strong for chart existence but requires specialist or manual transcription before precise dental notation is used in prose.

The death mask or bust is catalogued as PE-0027 and maps to EV-0007. Lawson made the cast, and the bust was marked as exhibit C.17. The record is high confidence for creation but low for current custody. Lawson's own limitations about formalin, embalming, and postmortem change restrict living-appearance inference from the object.

Hair evidence is catalogued as PE-0028 and maps to EV-0006. The repository has weak primary support for hair as a preserved exhibit in the acquired 1949 material, beyond body-description context. Later hair, death-mask, or exhumation claims require official sample-source, custody, and laboratory records before they can be used as established physical-evidence claims.

Biological specimens are catalogued as PE-0029 and map to EV-0017. The record supports that stomach contents, liver and muscle, urine, and blood were submitted to Cowan and analyzed in 1948. The present availability of those samples is unknown and likely unavailable, but the repository does not contain a primary retention or disposal record.

\section{Chain of Custody}

The chain-of-custody confidence matrix is the controlling register for physical-evidence weight. It records 34 custody entries. Twenty-nine are fragmentary, four are unknown, and one is not applicable. No catalogue item has a complete custody trail.

Fragmentary custody means that existence or observation may be well supported while later movement, storage, transfer, destruction, or current custodian is incomplete. This distinction is essential for the suitcase, clothing, slip, Rubaiyat material, thread, photographs, biological specimens, death mask, fingerprints, and dental chart.

Unknown custody means the item is catalogued but present handling cannot be reconstructed from the available repository records. PE-0023 sand, PE-0024 vegetation, PE-0028 hair, and PE-0033 cash fall into this class. These items remain inventory entries, but their evidentiary weight is limited.

\begin{table}[htbp]
\centering
\caption{Chain-of-custody confidence summary}
\label{tab:chain-of-custody-summary-ch6}
\begin{tabularx}{\textwidth}{>{\raggedright\arraybackslash}p{0.28\textwidth}r>{\raggedright\arraybackslash}X}
\toprule
Custody classification & Count & Meaning in this chapter \\
\midrule
Complete & 0 & No exhibit has an end-to-end custody trail in the repository. \\
Mostly complete & 0 & No exhibit has only minor custody gaps. \\
Fragmentary & 29 & Existence or observation is documented, but later transfer, storage, current location, or destruction is incomplete. \\
Unknown & 4 & Item-level custody cannot presently be reconstructed from repository evidence. \\
Lost & 0 & No item is classified as lost on primary-source proof alone. \\
Not applicable & 1 & Environmental context is not a recoverable object. \\
\bottomrule
\end{tabularx}
\end{table}

\section{Preservation History}

The preservation register contains 10 entries. Three are preserved locally as digital copies: the Rubaiyat code image, fingerprint chart, and dental chart. Two are located only as references: body photographs and suitcase photographs. The death mask or bust has unverified current status. The suitcase is unknown with destruction not primary-verified. Body clothing and the Tamam Shud slip have unknown current status. Biological specimens are likely unavailable, but that is not primary-verified.

The preservation register should therefore be read as an acquisition guide. It identifies what the project can presently inspect, what it can only reference, and what requires custodian confirmation before publication or interpretation. It also prevents public-domain or online image availability from being confused with exhibit custody.

\begin{table}[htbp]
\centering
\caption{Preservation timeline and status}
\label{tab:preservation-timeline-ch6}
\begin{tabularx}{\textwidth}{>{\raggedright\arraybackslash}p{0.14\textwidth}>{\raggedright\arraybackslash}p{0.22\textwidth}>{\raggedright\arraybackslash}p{0.24\textwidth}>{\raggedright\arraybackslash}X}
\toprule
ID & Item & Preserved form & Status \\
\midrule
PRH-0001 & Rubaiyat code image & Digital JPG reproduction & Preserved locally as digital copy; rights and original custody unclear. \\
PRH-0002 & Fingerprint chart & Digital PDF reproduction & Preserved locally as digital copy; original police or archive custodian unknown. \\
PRH-0003 & Dental chart & Digital PDF reproduction & Preserved locally as digital copy; original custody unknown. \\
PRH-0004 & Body photographs & Referenced in DOC-0001 & Located as reference only; image set and rights not established. \\
PRH-0005 & Suitcase photographs & Referenced in DOC-0001 & Located as reference only; image set and rights not established. \\
PRH-0006 & Death mask or bust & Physical object reported in testimony & Current status unverified in acquired primary archive. \\
PRH-0007 & Suitcase & Physical object or exhibit in 1949 & Current status unknown; destruction not primary-verified. \\
PRH-0008 & Body clothing & Physical exhibit in 1949 & Current status unknown. \\
PRH-0009 & Tamam Shud slip & Physical exhibit in 1949 & Current status unknown. \\
PRH-0010 & Biological specimens & 1948 specimens documented & Likely unavailable, but no primary retention or disposal record acquired. \\
\bottomrule
\end{tabularx}
\end{table}

\section{Physical Evidence Assessment}

The strongest physical-evidence records are the body location and position, body clothing and initial pocket items, rail and bus tickets, suitcase existence and police possession, Tamam Shud slip existence and fob-pocket association, fingerprint creation and circulation, dental chart creation, body photographs, and bust creation. These records are strong because they are mapped to witness testimony, inquest pages, exhibit references, or preserved reproductions.

The weakest physical-evidence records are those where the object may have existed but present custody, exact wording, or primary-source support is incomplete. These include cash, vegetation, hair as later physical evidence, sand as a retained sample, the current status of the death mask, the original associated Rubaiyat, and several OCR-dependent suitcase inventory details.

The principal missing evidence is not a single object but a set of custody and preservation records: original body photographs C.5--C.8, suitcase photographs C.10--C.14, full clothing photographs or measurements, original property inventory, original suitcase layout, original Rubaiyat provenance, the Tamam Shud slip's discovery log, thread and label retention records, biological-specimen retention or disposal records, and current custodian records for major exhibits.

\begingroup
\scriptsize
\setlength{\tabcolsep}{2pt}
\renewcommand{\arraystretch}{1.12}
\begin{longtable}{>{\raggedright\arraybackslash}p{0.10\textwidth}>{\raggedright\arraybackslash}p{0.13\textwidth}>{\raggedright\arraybackslash}p{0.32\textwidth}>{\raggedright\arraybackslash}p{0.15\textwidth}>{\raggedright\arraybackslash}p{0.20\textwidth}}
\caption{Complete physical evidence inventory}\label{tab:complete-physical-evidence-inventory}\\
\toprule
ID & Evidence ID & Description & Custody & Confidence \\
\midrule
\endfirsthead
\toprule
ID & Evidence ID & Description & Custody & Confidence \\
\midrule
\endhead
PE-0001 & EV-0005 & Body of unknown male found near Somerton Beach seawall. & Fragmentary & High initial movement; Moderate later history \\
PE-0002 & EV-0014 & Partly smoked cigarette near cheek, lapel, or collar area. & Fragmentary & Moderate \\
PE-0003 & EV-0005 & Clothing worn by deceased as an ensemble. & Fragmentary & High existence; Moderate details \\
PE-0004 & EV-0004 & Body tie and suitcase ties, including T. Keane tie reference. & Fragmentary & Moderate \\
PE-0005 & EV-0005 & Shoes worn by deceased. & Fragmentary & High condition observations \\
PE-0006 & EV-0005 & Socks worn by deceased. & Fragmentary & Moderate \\
PE-0007 & N/A & Removed or missing tags and labels. & Fragmentary & Moderate \\
PE-0008 & EV-0008 & Second-class Adelaide-to-Henley Beach railway ticket. & Fragmentary & High \\
PE-0009 & EV-0009 & Municipal Tramways Trust bus ticket. & Fragmentary & High \\
PE-0010 & EV-0012 & Army Club cigarette packet containing Kensitas cigarettes. & Fragmentary & Moderate \\
PE-0011 & EV-0013 & Bryant and May matches box. & Fragmentary & Moderate \\
PE-0012 & EV-0010 & Two combs or metal comb body property. & Fragmentary & Moderate \\
PE-0013 & EV-0011 & Juicy Fruit chewing gum packet. & Fragmentary & Moderate \\
PE-0014 & EV-0001 & Tamam Shud paper slip from fob pocket. & Fragmentary & High existence; Moderate custody \\
PE-0015 & EV-0002 & Rubaiyat copy or reference book associated with comparison. & Fragmentary & Moderate \\
PE-0016 & N/A & Rubaiyat code or markings digital image. & Fragmentary & Moderate digital; Low original custody \\
PE-0017 & EV-0003 & Unclaimed suitcase from Adelaide Railway Station. & Fragmentary & High inquest existence; Moderate custody \\
PE-0018 & N/A & Railway Station luggage ticket no. 52703. & Fragmentary & High \\
PE-0019 & EV-0003 & Suitcase clothing and textiles. & Fragmentary & Moderate \\
PE-0020 & N/A & Suitcase tools and personal articles. & Fragmentary & Moderate \\
PE-0021 & N/A & Six pencils or writing implements in suitcase. & Fragmentary & Moderate \\
PE-0022 & N/A & Orange, tan, or warm sepia thread. & Fragmentary & High testimony; Moderate physical evidence \\
PE-0023 & N/A & Sand observations at scene, body, clothing, or suitcase. & Unknown & Moderate \\
PE-0024 & N/A & Barley grass or seed material. & Unknown & Low \\
PE-0025 & EV-0015 & Fingerprints or fingerprint chart. & Fragmentary & High existence; Moderate custody \\
PE-0026 & EV-0016 & Dental chart and dental observations. & Fragmentary & High existence; Moderate details \\
PE-0027 & EV-0007 & Death mask, mould, cast, or bust. & Fragmentary & High creation; Low current custody \\
PE-0028 & EV-0006 & Hair as potential later evidence. & Unknown & Low \\
PE-0029 & EV-0017 & Stomach contents, liver, muscle, urine, and blood specimens. & Fragmentary & High 1948 submission; Low present preservation \\
PE-0030 & EV-0003 & Suitcase photographs C.10--C.14. & Fragmentary & Moderate \\
PE-0031 & EV-0005 & Body photographs C.5--C.8. & Fragmentary & High creation; Moderate custody \\
PE-0032 & N/A & Scene environment and public beach conditions. & N/A & Moderate \\
PE-0033 & N/A & 2s 6d cash entry from OCR-dependent inventory. & Unknown & Low \\
PE-0034 & N/A & Stencil materials and related suitcase tools. & Fragmentary & Moderate \\
\bottomrule
\end{longtable}
\endgroup

\begin{table}[htbp]
\centering
\caption{Known missing or unlocated evidence}
\label{tab:known-missing-evidence-ch6}
\begin{tabularx}{\textwidth}{>{\raggedright\arraybackslash}p{0.16\textwidth}>{\raggedright\arraybackslash}p{0.28\textwidth}>{\raggedright\arraybackslash}X}
\toprule
Record & Evidence area & Missing or unresolved material \\
\midrule
PRH-0004 & Body photographs & Original C.5--C.8 image set, custody, and publication rights. \\
PRH-0005 & Suitcase photographs & Original C.10--C.14 image set, custody, and publication rights. \\
PRH-0006 & Death mask or bust & Current custodian, accession number, condition, and access status. \\
PRH-0007 & Suitcase & Retention or destruction authority and present status. \\
PRH-0008 & Body clothing & Present location, preservation, and item-level measurements or images. \\
PRH-0009 & Tamam Shud slip & Present location, best image, and custody file. \\
PRH-0010 & Biological specimens & Retention or disposal record and any surviving laboratory notes. \\
\bottomrule
\end{tabularx}
\end{table}

\begin{table}[htbp]
\centering
\caption{Evidence uncertainty register}
\label{tab:evidence-uncertainty-register-ch6}
\begin{tabularx}{\textwidth}{>{\raggedright\arraybackslash}p{0.14\textwidth}>{\raggedright\arraybackslash}p{0.18\textwidth}>{\raggedright\arraybackslash}X>{\raggedright\arraybackslash}p{0.14\textwidth}}
\toprule
ID & Evidence & Question & Priority \\
\midrule
PEU-0001 & PE-0014 & When exactly was the Tamam Shud slip found and by whom? & Critical \\
PEU-0002 & PE-0017 & Who deposited the suitcase, and does the ticket prove association with the deceased? & Critical \\
PEU-0003 & PE-0007 & Were labels deliberately removed before death? & High \\
PEU-0004 & PE-0002 & Was the partly smoked cigarette collected and examined? & High \\
PEU-0005 & PE-0022 & Was the thread physically retained or only observed? & High \\
PEU-0006 & PE-0031 & Where are the body photographs C.5--C.8? & High \\
PEU-0007 & PE-0030 & Where are suitcase photographs C.10--C.14? & High \\
PEU-0008 & PE-0029 & Were biological samples retained? & Critical \\
PEU-0009 & PE-0027 & What is the current status or custodian of the death mask or bust? & High \\
PEU-0010 & PE-0015 & What happened to the associated Rubaiyat copy? & Critical \\
\bottomrule
\end{tabularx}
\end{table}
